% ACM artifact appendix (standalone).  Build: ./build.sh artifact
%
% It is deliberately a separate document: ACM artifact evaluation is a separate
% submission from the paper, and the paper appendix must not be the only place
% where the reproduction commands live.
\documentclass[acmsmall,10pt,nonacm]{acmart}
\usepackage{booktabs}
\usepackage{array}
\usepackage{lmodern}
\settopmatter{printfolios=true,printacmref=false}
\newcommand{\code}[1]{\texttt{\detokenize{#1}}}
\setlength{\emergencystretch}{2em}
\begin{document}

\title{Artifact Appendix: Bounded Reclamation}
\subtitle{Compile-time tiering, pause and memory contracts, and a roofline for reclamation policy}
\author{Anonymous Author(s)}
\affiliation{\institution{Paper under double-blind review}\country{Review copy}}

\begin{abstract}
This artifact is the extC compiler together with the harness that produces every number in the paper:
the four-shape cross-language matrix (eleven arms of C++, Rust, Go, Java and JavaScript), the pause
harnesses, the ten-variant mechanism ablation, the boundedness runs, the tier census, the Lean~4
development, and the gates that hold the runtime to its budgets.  Everything is built from source with
\code{make}; no pre-built binary is required and no external data set is used.
\end{abstract}

\maketitle

\section{Artifact check-list (meta-information)}
\begin{table}[h]
\small
\begin{tabular}{@{}p{0.26\linewidth}p{0.68\linewidth}@{}}
\toprule
\textbf{Item} & \textbf{Value} \\
\midrule
Program & extC compiler (C, \code{src/}, built with \code{make}); generated C links only against libc \\
Runtime library & none: the reclamation runtime is emitted into each translated unit (\code{src/back/reclaim.c}) \\
Proofs & Lean~4, thirteen modules, \code{proofs/}; optional for the experiments \\
Tests & 216-case C driver (\code{tests/reclaim}), emitted-code gates (\code{tools/}), whole-program instruction budget \\
Data set & none external; every input is generated from a seed at run time \\
Hardware & x86-64, four or more cores, 8\,GiB free; a single pinned core is used for timing \\
Run-time environment & Linux, \code{gcc}, \code{make}, \code{python3} with \code{matplotlib} \\
Optional tool-chains & \code{rustc}, \code{go}, \code{javac}/\code{java}, \code{node} for competitor arms; \code{valgrind} for the instruction gate \\
Metrics & wall clock, peak RSS, pause samples, retired instructions (callgrind), emitted tier \\
Output & \code{.tsv} matrices, \code{.md} reports, \code{.png}/\code{.pdf} figures per script \\
Experiments & four shapes $\times$ eleven arms (7 interleaved reps); ten-variant ablation; boundedness; census \\
Publicly available & yes, to be hosted at an anonymized URL for review and archived with a DOI at publication \\
Code license & Apache-2.0 (\code{LICENSE}) \\
\bottomrule
\end{tabular}
\caption{Meta-information for the artifact evaluation.}
\end{table}

\section{Description}
The repository is self-contained.  \code{make} builds \code{build/extc}.  The four shape programs and
the eleven comparison arms live in \code{bench/gc/v2}, their pause harnesses in
\code{bench/gc/v2/pause}, the driver suite in \code{tests/reclaim}, the Lean development in
\code{proofs}, the paper sources in \code{paper/} (with \code{paper/data/*.tsv} holding the checked-in
numbers), and the scripts that generate every table and figure in \code{paper/figures}.  The runtime
constants that the ablation moves (\code{EXTC\_RC\_SLAB\_BITS}, \code{EXTC\_RC\_SCAN\_CHUNK},
\code{EXTC\_RC\_SEG\_DEADLINE\_US}, \code{EXTC\_RC\_SWEEP\_FINISH\_MAX}, and the rest) are
\code{\#ifndef}-guarded, which is what makes the one-constant-at-a-time ablation possible without
touching the compiler.

\section{Creation}
\begin{itemize}
\item \textbf{How to build:} \code{make} at the repository root, then \code{make check} (or
  \code{./tests/run.sh}) for the whole test suite.  A minimal build needs only \code{gcc}.
\item \textbf{How to compile one program:} \code{./build/extc -o out.c prog.extc} emits C;
  \code{cc -O2 out.c -o prog} links it.  Adding \code{EXTC\_RECLAIM=0} at translation time removes the
  runtime entirely and is the ``off'' arm of every ablation.
\item \textbf{How to check the proofs:} \code{cd proofs \&\& bash verify.sh} checks the Lean~4
  development.  The modules must be compiled in dependency order first (the shipped script assumes the
  \code{.olean} files exist, and they are not in the repository); the audit and the dependency order
  are in \code{proofs/README.md}.
\item \textbf{No network access} is needed once the repository is unpacked.
\end{itemize}

\section{Getting started (verified path)}
\begin{enumerate}
\item \code{make}, then \code{./build/extc -o /tmp/hello.c examples/new.extc \&\& cc -O2 /tmp/hello.c -o /tmp/hello}.
\item \code{tests/reclaim/run.sh} (expect 216 passing cases; the suite also runs under
  \code{-fsanitize=address}, \code{-fsanitize=undefined} and with the runtime out of line).
\item \code{cd bench/gc/v2 \&\& B=/tmp/ar EXTC=../../build/extc bash run.sh align} builds every arm,
  then \code{B=/tmp/ar TSV=/tmp/ar/matrix.tsv REPS=7 PIN=6 bash run.sh bench} runs the matrix and
  \code{python3 report.py /tmp/ar/matrix.tsv REPORT.md matrix.png pause\_v2.tsv} renders it.
\item \code{python3 paper/figures/gen\_tables.py} and \code{gen\_roofline.py} regenerate every table and
  figure of the paper from \code{paper/data/*.tsv}.
\end{enumerate}

\section{Evaluation and expected results}
\begin{itemize}
\item \textbf{Correctness.} Every arm of the matrix prints the same two lines of output; the harness
  refuses to time an arm that does not, so a run that completes is already a differential test.
\item \textbf{Throughput.} On a pinned core, the four extC numbers are approximately 39, 16, 117 and
  15\,ms for \code{promote}, \code{sizeratio}, \code{chase} and \code{cycleratio} at the paper's sizes;
  absolute values move with the machine and between sessions by up to 13\%, so only ratios formed inside
  one interleaved run should be compared.
\item \textbf{Memory.} Peak RSS is read with \code{/usr/bin/time -f \%M}; the boundedness run keeps
  3,420\,kB at 125\,000, 1\,000\,000 and 8\,000\,000 rounds with reclamation on, and grows by
  37$\times$ with it off.
\item \textbf{Pause.} With \code{EXTC\_RC\_PAUSE=1 EXTC\_RC\_SLOWSTEP=0.001} the runtime logs one line
  per step; no step on any shape exceeds one millisecond, and the 99.9th percentile is 10--58\,$\mu$s.
\item \textbf{Budget gate.} \code{tests/reclaim-perf/gate.sh} reruns one shape under callgrind and fails
  above the frozen instruction budget; \code{tests/reclaim-gaps/gate-release.sh} holds the emitted-code
  properties; \code{tools/check\_tree.py --gate} holds the tree hygiene.
\end{itemize}

\section{Methodology and pitfalls}
Runs are interleaved round-robin with reversed order on alternate rounds, one core pinned with
\code{taskset}, and every arm warmed once before its first timed round.  Do not build or run anything
else on the machine during a measurement round: the pause instrumentation charges preempted time to the
step that was preempted, and a single worst step is a report about the machine, not about the runtime.
Report medians of the per-round minimum, and never compare absolute milliseconds across sessions.

\section{Anonymization status (review)}
This appendix and the paper are built from the same source and contain no author names, affiliations,
acknowledgements or personal URLs.  The artifact repository is to be uploaded to an anonymized host for
review; the camera-ready version substitutes \code{authors.tex} and the archival DOI.  The proof scripts
shell out to \code{lean} and (optionally) \code{z3}; neither is needed for the experiments.

\end{document}
